\documentclass[11pt,a4paper]{article}
\usepackage[T1]{fontenc}
\usepackage{lmodern,amsmath,amssymb,amsthm,mathtools,booktabs,hyperref}
\usepackage[margin=27mm]{geometry}
\setlength{\emergencystretch}{2em}
\newtheorem{theorem}{Theorem}[section]
\newtheorem{proposition}[theorem]{Proposition}
\newtheorem{lemma}[theorem]{Lemma}
\newtheorem{corollary}[theorem]{Corollary}
\theoremstyle{remark}\newtheorem{remark}[theorem]{Remark}
\newcommand{\Q}{\mathbb Q}\newcommand{\Z}{\mathbb Z}
\newcommand{\Ga}{\mathbb G_a}\newcommand{\Aa}{\mathbb A}
\DeclareMathOperator{\Res}{Res}
\title{Structural constraints on reductions of elliptic witnesses}
\author{}\date{}
\begin{document}\maketitle
\begin{abstract}
We examine algebraic reductions of a rational-source construction
whose integer labels are carried by multiples of a point over a
quadratic number field.  We prove dimension bounds for recovering
arbitrary tuples of such points, the Zariski density of the graph
of the integer label, and an obstruction to replacing the paired
conjugate reductions by one rational elliptic curve.  A related
two-curve obstruction applies when a curve correspondence is used
to enforce a common integer index.  We also identify the exact
squareclass restriction on a shared pole parameter and give an
explicit counterexample to a parity shortcut in scalar height
arguments.  These statements concern the specified strategies;
they are not lower bounds for arbitrary Diophantine definitions.
\end{abstract}

\section{Point labels and finite algebraic presentations}

Let $K=\Q(\sqrt2)$ and
\[
 D=11-6\sqrt2,\qquad E:y^2=x^3-Dx,
\]
\[
 P=\left(\frac{57-28\sqrt2}{8},
                \frac{-660+417\sqrt2}{32}\right).
\]
Let $A=\Res_{K/\Q}E$ be the Weil restriction, and identify
$E(K)$ with $A(\Q)$.  The elliptic-chart argument in~\cite{Charts}
proves that $A$ is simple over $\Q$; in particular its proof
distinguishes the two conjugate curves by their good reductions
at $17$.  The point $P$ has infinite order.  Write
$\Gamma=\Z P\subset A(\Q)$.

\begin{lemma}\label{lem:dense}
The subgroup $\Gamma$ is Zariski dense in $A$, and
$\Gamma^r$ is Zariski dense in $A^r$ for each positive integer $r$.
\end{lemma}
\begin{proof}
The Zariski closure of a subgroup of an algebraic group is an
algebraic subgroup.  Its identity component is defined over $\Q$.
Since $P$ has infinite order, this component has positive dimension;
simplicity forces it to be $A$.  Density of the Cartesian powers
follows successively by restricting a closed subset to fibers in
one factor.
\end{proof}

\begin{theorem}\label{thm:dimension}
Suppose finitely many algebraic varieties $Y_j$ over $\Q$ admit
finite morphisms to locally closed subvarieties of $\Aa^b$, and
morphisms $F_j:Y_j\to A^r$.  If
\[
            \Gamma^r\subset\bigcup_jF_j(Y_j(\Q)),
\]
then $b\ge2r$.
\end{theorem}
\begin{proof}
Each $Y_j$ has dimension at most $b$.  The finite union of the
closures of the images of the $F_j$ contains $\Gamma^r$, and
therefore equals $A^r$ by Lemma~\ref{lem:dense}.  Since $A^r$ is
irreducible, one of these image closures is all of $A^r$.
Consequently $2r=\dim A^r\le\dim Y_j\le b$.
\end{proof}

The theorem applies to algebraic presentations in which $b$
rational coordinates are retained and all further coordinates are
finite algebraic choices.  Localizing by a nonzero guard or taking
finitely many charts does not change the argument.  In particular,
recovering arbitrary tuples of nine point labels in this way requires
at least eighteen retained rational parameters.  If an input
coordinate and one final finite-algebra coordinate are counted
separately, fewer than ten total coordinates permit at most three
independent point labels in this generic recovery model.

This conclusion is deliberately limited.  Completeness of a
Diophantine construction need not require arbitrary tuples of
labels, and its arithmetic witnesses can satisfy useful dependencies.
The theorem does not exclude such dependencies or a different
integer-label mechanism.

\section{The integer index is not an algebraic coordinate}

\begin{theorem}\label{thm:index-graph}
Let $B$ be a positive-dimensional abelian variety over a field of
characteristic zero, and let $Q\in B$ generate a Zariski-dense
cyclic subgroup.  Then
\[
             \{(n,[n]Q):n\in\Z\}
\]
is Zariski dense in $\Ga\times B$.
The same assertion holds with $n$ restricted to any nonconstant
arithmetic progression.
\end{theorem}
\begin{proof}
The displayed set is the cyclic subgroup generated by $(1,Q)$.
Let $G$ be its Zariski closure, and pass to the identity component
$G^0$.  Both projections of $G^0$ are surjective: the images of
$G$ are dense subgroups of the respective factors, and the finite
component group cannot account for a positive-dimensional connected
quotient.

The kernel of $G^0\to B$ is an algebraic subgroup of $\Ga$, hence
is either trivial or all of $\Ga$ in characteristic zero.  In the
first case the surjective homomorphism has finite kernel and is
an isogeny of connected algebraic groups.  Such a homomorphism is
finite, so $G^0$ is proper because $B$ is proper.  Every morphism
from a connected proper variety to the affine line is constant,
contradicting the surjective projection to $\Ga$.  Thus the kernel
contains $\Ga$.  Surjectivity to $B$ then gives
$G^0=\Ga\times B$.

For an arithmetic progression $n=a+dm$, with $d\ne0$, the closure
of $\Z[d]Q$ is again all of $B$: a proper closed subgroup containing
it would have finitely many translates containing $\Z Q$.
The preceding argument applied to $(d,[d]Q)$ proves the assertion,
after translation by $(a,[a]Q)$.
\end{proof}

\begin{corollary}\label{cor:no-index}
On a nonempty affine open subset of $A$, no nonzero polynomial
relation with coefficients in its coordinate ring holds between
the integer $n$ and the point $[n]P$ for all integers $n$ for which
the point belongs to that open set.  In particular, an algebraic
finite extension of the point coordinates cannot recover every
integer label.
\end{corollary}
\begin{proof}
Any such relation would cut out a proper closed subset of the
product of the open set with $\Aa^1$, contrary to
Theorem~\ref{thm:index-graph}.
\end{proof}

For example, the Lucas relation
\[
 t^2-a(a-2)tb+(a-2)^2b^2-1=0
\]
is monic in $t$.  If $a,b$ are already integers and $t$ is rational,
the rational-root theorem gives $t\in\Z$.  In the elliptic source,
however, $a,b,t$ are integer indices of points, rather than
available rational coordinates.  Replacing a point witness by this
monic argument requires a separate way to make the relevant
indices available.  Corollary~\ref{cor:no-index} rules out doing so
by a uniform finite algebraic recovery of all labels.

\section{The paired-place mechanism over a rational curve}

Let $E_0/\Q$ be an elliptic curve, let $P_0\in E_0(\Q)$ be
non-torsion, and let $K/\Q$ be quadratic.  At a rational prime
$p$ of good reduction which splits in $K$, the two residue fields
are canonically $\mathbb F_p$.  A point belongs to the formal
group at one of the two places exactly when its reduction there
is the identity.

\begin{proposition}\label{prop:rational-curve}
If $H=[h]P_0$ is formal at one place above $p$, and $H-P_0$ is
formal at the other, then $P_0$ reduces to the identity modulo $p$.
Only finitely many good rational primes can satisfy this conclusion
for a fixed affine point $P_0$.
\end{proposition}
\begin{proof}
The reductions of $P_0$ at the two places are the same rational
point $\overline P_0\in E_0(\mathbb F_p)$.  The two requirements give
\[
              [h]\overline P_0=O,
              \qquad[h-1]\overline P_0=O.
\]
Subtracting yields $\overline P_0=O$.
For a fixed affine point this occurs only at primes dividing
the denominators of its coordinates in a fixed integral model,
together with the finitely many excluded primes.
\end{proof}

Thus a rational elliptic curve cannot replace the quadratic-field
curve while retaining this paired-place test at arbitrarily many
contact primes.  The proposition concerns this mechanism, not
arbitrary constructions on rational elliptic curves.

\subsection{Two rational curves with a common index}

\begin{proposition}\label{prop:two-curves}
Let $E_1,E_2$ be elliptic curves over $\Q$ and let
$P_i\in E_i(\Q)$ be non-torsion.
If an algebraic curve in $E_1\times E_2$ contains
\[
                  \{([n]P_1,[n]P_2):n\in\Z\},
\]
then $E_1,E_2$ are isogenous over $\Q$.  There are an isogeny
$\phi:E_1\to E_2$, a nonzero integer $a$, and a torsion point
$T$ such that
\begin{equation}\label{eq:isogeny-relation}
                       \phi(P_1)=aP_2+T.
\end{equation}
For such a relation, coprimality of the two reduction orders can
occur at only finitely many good rational primes.
\end{proposition}
\begin{proof}
The Zariski closure of the displayed cyclic subgroup is an
algebraic subgroup of dimension at most one.  Its identity
component $G$ has dimension one because each $P_i$ is non-torsion.
The two projections $G\to E_i$ are isogenies.  Let the first
projection have degree $a>0$, and compose its dual isogeny with
the second projection to obtain $\phi:E_1\to E_2$.
Some nonzero integer multiple of $(P_1,P_2)$ lies in $G$.
The defining property of the dual isogeny then shows that
$\phi(P_1)-aP_2$ is torsion, proving~\eqref{eq:isogeny-relation}.

Let $m$ be the order of $T$, let $d=\deg\phi$, and let $k_i$
be the reduction order of $P_i$ at a good prime where the
isogeny and the relation reduce.  Multiplying the relation gives
\[
                       k_2\mid amk_1.
\]
Also $[mk_2]P_1$ lies in the kernel of the reduced isogeny,
which is killed by $d$, and hence
\[
                       k_1\mid dmk_2.
\]
If $\gcd(k_1,k_2)=1$, these imply $k_1\mid dm$ and $k_2\mid am$.
Thus both orders are bounded independently of the prime.
For a fixed non-torsion point, bounded reduction order occurs
only at finitely many primes: each such prime makes one of
finitely many fixed nonzero multiples reduce to the identity.
\end{proof}

In particular, two nonisogenous rational curves may have the
needed distinct reduction orders, but their common-index pairs
cannot be recovered from a curve correspondence.  When a curve
correspondence does exist, the isogeny relation prevents
arbitrarily many coprime-order reductions.  A two-curve method
would therefore need a different way to enforce or replace
the common integer index.

\section{Sharing a pole parameter}

The first equation of each pole certificate is
\[
                  p_i^2=B_i(1-e_i)(1+3e_i),\qquad B_i\ne0.
\]

\begin{proposition}\label{prop:shared-e}
If two such equations share $e_i=e$ and $p_1p_2\ne0$, then
$B_1/B_2$ is a rational square.  If one of $p_1,p_2$ is zero,
then both are zero and $e\in\{1,-1/3\}$.
\end{proposition}
\begin{proof}
In the nonzero case division gives $B_1/B_2=(p_1/p_2)^2$.
Otherwise $(1-e)(1+3e)=0$, giving exactly the stated alternatives.
\end{proof}

The single rational chart used in the source is
\[
 e_i=\frac{\tau_i^2-B_i}{\tau_i^2+3B_i},\qquad
 p_i=-\frac{4B_i\tau_i}{\tau_i^2+3B_i}.
\]
Sharing $\tau_i$ is a different proposal: it need not share $e_i$,
so Proposition~\ref{prop:shared-e} does not apply.  It does,
however, impose a simultaneous rational-point requirement on
the remaining elliptic quartics.  The inherited realization
theorem chooses a fresh twist parameter for each prepared pole
input.  It does not assert that an arbitrary fixed or shared
chart parameter is available.

The degenerate value $e=-1/3$ is allowed by this chart, with
$\tau=0$.  Consequently the squareclass obstruction alone does
not exclude fixing this value.  Completeness of that specialization
would require a new rational-point theorem for its remaining
quartic; no such completeness conclusion is made here.

\section{Odd labels can be formal at both conjugate places}

\begin{proposition}\label{prop:odd-example}
For the curve and point of Section~1, both conjugate reductions
of $P$ at the split prime $23$ have order three.  Thus odd
integer labels can place their points in the formal groups
at both conjugate places simultaneously.
\end{proposition}
\begin{proof}
The two images of $\sqrt2$ in $\mathbb F_{23}$ are $5,18$.
The reduced data are
\[
\begin{array}{c|c|c|c}
 \sqrt2&D&P&[2]P\\\hline
 5&4&(4,5)&(4,18)\\
 18&18&(16,17)&(16,6).
\end{array}
\]
The last column is $-P$ in each row, and neither point is the
identity.  Hence both orders are three.  Every odd multiple
of three therefore reduces to the identity at both places.
\end{proof}

The exact checker also gives both orders equal to $25$ at
$p=113\equiv1\pmod8$.  Its complete enumeration below $3000$
finds $209$ split good primes, of which $72$ have odd orders
at both conjugate places.  These computations refute the
shortcut that parity of the index singles out one formal
embedding.  They do not refute any of the proved aggregate
height comparisons.

\section{Rational injections and finite fibers}

Pasten constructs a regular injection on rational points from
a suitable affine surface $V\times V$, with $V$ an open part
of a positive-rank elliptic curve, to the affine line~\cite{Pasten}.
Such a rational-point injection need not have finite geometric
fibers.  Indeed any dominant morphism from a surface to a curve
has positive-dimensional generic fiber.  It therefore does not
by itself provide a finite-algebra presentation with a single
retained coordinate.

Likewise, the polynomial injections of two rational variables
studied by Poonen~\cite{Poonen} concern injectivity on rational
points.  The stated existence theorem assumes a rational-point
nondensity hypothesis.  Neither injectivity nor this additional
hypothesis supplies the finite geometric fibers needed by the
finite-root elimination theorem.  An application to witness
compression must provide that separate step.

\begin{thebibliography}{9}
\bibitem{Charts}
\emph{An obstruction to one-parameter presentations of an elliptic point},
project paper, version 7, 7 November 2026.
\bibitem{Source}
\emph{A defining polynomial in twenty-seven variables:
an existence theorem with a uniform degree bound},
project paper, version 12, 9 November 2026.
\bibitem{Pasten}
H.~Pasten, \emph{Bivariate polynomial injections and elliptic curves},
arXiv:1811.02416, \url{https://arxiv.org/abs/1811.02416}.
\bibitem{Poonen}
B.~Poonen, \emph{Multivariable polynomial injections on rational numbers},
Acta Arith. \textbf{145} (2010), 123--127,
\url{https://math.mit.edu/~poonen/papers/QQ.pdf}.
\end{thebibliography}
\end{document}
